% =============================================================================
%  AWL: A Declarative Agent Workflow Language with Big-Step Operational
%       Semantics -- Course Report
%  Authors: Ankush Chandra, Gautam Panigrahi, Kunal Elchipuram
%  SJSU -- CS Dept, May 2026
% =============================================================================
\documentclass[10pt]{article}

% ---- Layout ------------------------------------------------------------------
\usepackage[margin=0.85in, top=1in, bottom=1in, columnsep=0.3in]{geometry}
\usepackage{multicol}

% ---- Math / Semantics --------------------------------------------------------
\usepackage{amsmath, amssymb, amsthm}
\usepackage{mathpartir}          % \inferrule, mathpar environment
\usepackage{stmaryrd}            % \llbracket, \rrbracket

% ---- Code listings -----------------------------------------------------------
\usepackage{listings}
\usepackage{xcolor}

\definecolor{lbg}{rgb}{0.97,0.97,0.97}
\definecolor{lcmt}{rgb}{0.25,0.50,0.25}
\definecolor{lkw}{rgb}{0.00,0.00,0.60}
\definecolor{lstr}{rgb}{0.50,0.00,0.50}
\definecolor{lframe}{rgb}{0.75,0.75,0.75}

\lstdefinestyle{hs}{
  backgroundcolor  = \color{lbg},
  commentstyle     = \color{lcmt}\itshape,
  keywordstyle     = \color{lkw}\bfseries,
  stringstyle      = \color{lstr},
  basicstyle       = \ttfamily\scriptsize,
  breaklines       = true,
  keepspaces       = true,
  showstringspaces = false,
  tabsize          = 2,
  language         = Haskell,
  morekeywords     = {data,where,case,of,do,let,in,return,type,
                      newtype,deriving,module,import,instance,class,
                      Map,IO,String,Bool,Int,Double,Integer,Maybe},
  frame            = single,
  rulecolor        = \color{lframe},
  framesep         = 2pt,
  xleftmargin      = 4pt,
  xrightmargin     = 4pt,
  belowskip        = 4pt,
  aboveskip        = 4pt
}

\lstdefinelanguage{AWL}{
  keywords         = {config,agent,from,let,if,then,else,fail,retry,
                      try,catch,print,FixedAgent,CustomAI,prompt,model,
                      true,false,null,python,http,llm,mock},
  sensitive        = true,
  comment          = [l]{//},
  morecomment      = [s]{/*}{*/},
  morestring       = [b]"
}
\lstdefinestyle{awl}{
  backgroundcolor  = \color{lbg},
  commentstyle     = \color{lcmt}\itshape,
  keywordstyle     = \color{lkw}\bfseries,
  stringstyle      = \color{lstr},
  basicstyle       = \ttfamily\scriptsize,
  breaklines       = true,
  keepspaces       = true,
  showstringspaces = false,
  tabsize          = 2,
  language         = AWL,
  frame            = single,
  rulecolor        = \color{lframe},
  framesep         = 2pt,
  xleftmargin      = 4pt,
  xrightmargin     = 4pt,
  belowskip        = 4pt,
  aboveskip        = 4pt
}

% ---- Typography --------------------------------------------------------------
\usepackage[T1]{fontenc}
\usepackage{microtype}
\usepackage{parskip}
\setlength{\parskip}{4pt}
\setlength{\parindent}{0pt}

% ---- Tables / Floats ---------------------------------------------------------
\usepackage{booktabs}
\usepackage{caption}
\captionsetup{font=small, labelfont=bf, skip=4pt}

% ---- Hyperlinks --------------------------------------------------------------
\usepackage[hidelinks]{hyperref}

% ---- List spacing ------------------------------------------------------------
\usepackage{enumitem}
\setlist{itemsep=2pt, topsep=3pt, parsep=0pt}

% ---- Section heading sizes ---------------------------------------------------
\usepackage{titlesec}
\titleformat{\section}{\normalsize\bfseries}{\thesection.}{0.5em}{}
\titleformat{\subsection}{\small\bfseries}{\thesubsection.}{0.5em}{}
\titleformat{\subsubsection}{\small\itshape}{\thesubsubsection.}{0.4em}{}
\titlespacing*{\section}{0pt}{8pt}{4pt}
\titlespacing*{\subsection}{0pt}{5pt}{2pt}

% ---- Semantic macros ---------------------------------------------------------
\newcommand{\A}{\mathcal{A}}
\newcommand{\J}{\A;\sigma;\Gamma \vdash}          % standard judgment lhs
\newcommand{\Jp}{\A;\sigma';\Gamma' \vdash}        % primed judgment lhs
\newcommand{\err}[1]{\mathrm{err}(#1)}
\newcommand{\sem}[1]{\llbracket #1 \rrbracket}
\newcommand{\Dom}{\mathrm{dom}}
\newcommand{\backend}{\mathit{backend}}
\newcommand{\fixed}{\mathit{fixed}}
\newcommand{\custom}{\mathit{custom}}
\newcommand{\runfixed}{\mathit{run\_fixed}}
\newcommand{\llmcall}{\mathit{llm\_call}}
\newcommand{\apply}{\mathit{apply}}

% =============================================================================
\begin{document}

% ---- Title (single column) --------------------------------------------------
\begin{center}
  {\Large\bfseries AWL: A Declarative Agent Workflow Language\\[4pt]
   with Big-Step Operational Semantics}\\[10pt]
  {\normalsize
    Ankush Chandra \quad Gautam Panigrahi \quad Kunal Elchipuram}\\[3pt]
  {\small Department of Computer Science, San Jose State University}\\[3pt]
  {\small May 2026}
\end{center}

% ---- Abstract (single column) -----------------------------------------------
\begin{center}
\begin{minipage}{0.88\linewidth}
\small\textbf{Abstract.}
We present AWL (Agent Workflow Language), a statically-structured
domain-specific language for composing multi-agent AI workflows.  AWL provides
first-class constructs for declaring agents backed by Python modules, HTTP
endpoints, large language models, or deterministic fixed-agent roles, and for
composing them with sequential bindings, conditionals, retries, and structured
error recovery.  We formalise the language with a big-step operational semantics
using three environments---variable~$\sigma$, agent~$\A$, and
configuration~$\Gamma$---and define evaluation judgements for both expressions
and statements.  The language is implemented in Haskell using Parsec for
parsing and a monadic big-step interpreter.  A suite of 49 integration tests
validates the implementation against the formal specification.
\end{minipage}
\end{center}

\vspace{6pt}
\hrule
\vspace{6pt}

% ---- Body (two columns) -----------------------------------------------------
\begin{multicols}{2}

% =============================================================================
\section{Introduction}
\label{sec:intro}

The proliferation of large language model (LLM) APIs has enabled a new paradigm
of \emph{agentic} computation, in which complex tasks are decomposed across
multiple co-operating AI agents.  Frameworks such as LangChain~\cite{langchain},
AutoGen~\cite{autogen}, and n8n~\cite{n8n} provide workflow orchestration
capabilities, but they do so through library APIs or visual editors that lack a
formal specification.

Without a formal semantics it is difficult to reason about error propagation,
determinism of agent dispatch, the scope of variable bindings, or the precise
conditions under which a workflow may fail or recover.  Developers must rely on
documentation and runtime experimentation to understand pipeline behaviour.

\subsection{Motivation and Problem Statement}

We identify three concrete gaps in the current state of the art:

\begin{enumerate}
  \item \textbf{No formal model.}  Existing orchestration tools lack an
        operational semantics, making formal reasoning about correctness
        impossible.
  \item \textbf{Implicit error handling.}  Error propagation rules are
        ad hoc and backend-specific rather than defined at the language level.
  \item \textbf{Agents as library objects.}  Current frameworks treat agents as
        library instances rather than named, first-class language entities with
        declared backends and arities.
\end{enumerate}

\subsection{Contributions}

This paper makes the following contributions:

\begin{itemize}
  \item A declarative surface syntax for multi-agent workflow programs
        (\S\ref{sec:syntax}).
  \item A big-step operational semantics with three-environment judgements
        covering expressions, statements, error propagation, retry, and
        try/catch (\S\ref{sec:semantics}).
  \item A complete Haskell implementation: Parsec lexer/parser, monadic
        big-step interpreter, deterministic fixed-agent stubs, and optional
        live LLM backend integration (\S\ref{sec:implementation}).
  \item A 49-case integration test suite (\S\ref{sec:testing}).
\end{itemize}

% =============================================================================
\section{Language Overview}
\label{sec:syntax}

An AWL program is a sequence of statements evaluated top-to-bottom.
The language is dynamically typed: runtime values are strings, numbers,
booleans, records, lists, null, or error payloads.  There is no static
type-checker; type errors surface as runtime failures.

\subsection{Abstract Syntax}

The abstract syntax is defined in \texttt{Syntax.hs} and the formal
specification \texttt{agent\_workflow\_language.txt}.

\medskip
\noindent\textbf{Expressions}
\[
\begin{array}{rcl}
e & ::= & c \;\mid\; x \;\mid\; e.f \;\mid\; A(e_1,\ldots,e_n) \\
  & \mid & e_1 \oplus e_2 \;\mid\; \{f_i{=}e_i\} \;\mid\; e \rhd A
\end{array}
\]
where $c$ is a constant (string, number, boolean, or null), $x$ a variable,
$e.f$ field projection, $A(\ldots)$ an agent call, $e_1 \oplus e_2$ a
binary operation, $\{f_i{=}e_i\}$ a record literal, and $e \rhd A$
a pipeline expression.

\medskip
\noindent\textbf{Statements}
\[
\begin{array}{rcl}
s & ::= & \texttt{config}\;\{c_i{=}e_i\} \\
  & \mid & \texttt{agent}\;A\;\texttt{from}\;b \\
  & \mid & \texttt{agent}\;A = \texttt{FixedAgent}(k) \\
  & \mid & \texttt{agent}\;A = \texttt{CustomAI}(\textit{prompt}{=}e,\textit{model}{=}m)\\
  & \mid & \texttt{let}\;x = e \\
  & \mid & \texttt{if}\;e\;\texttt{then}\;s_1\;\texttt{else}\;s_2 \\
  & \mid & s_1\,;\,s_2 \;\mid\; \texttt{fail}\;e \\
  & \mid & \texttt{retry}\;n\;s \\
  & \mid & \texttt{try}\;s_1\;\texttt{catch}\;x \Rightarrow s_2
\end{array}
\]

\noindent\textbf{Backends and Fixed Kinds}
\[
\begin{array}{rcl}
b & ::= & \texttt{python:}"f" \mid \texttt{http:}"u"
    \mid \texttt{llm:}"m" \mid \texttt{mock:}"v"\\[3pt]
k & ::= & \textit{Planner} \mid \textit{TaskSplitter}
    \mid \textit{Extractor} \mid \textit{Critic} \\
  &     & \mid\;\textit{Writer} \mid \textit{Summarizer}
    \mid \textit{Validator} \mid \textit{Guardrail} \\
  &     & \mid\;\textit{Router} \mid \textit{Merger} \mid \textit{Ranker}
\end{array}
\]

\subsection{Values}
\[
v \;::=\; c \;\mid\; \{f_1{=}v_1,\ldots,f_n{=}v_n\}
         \;\mid\; \err{v}
\]
Error values $\err{v}$ represent runtime failures carrying a payload $v$
(typically a descriptive string).  In the Haskell implementation they are
represented as the type constructor \texttt{EErr Value} rather than being
embedded inside the value domain.

\subsection{Runtime Environments}
\[
\begin{array}{rcl}
\sigma & : & \textit{Var} \rightharpoonup \textit{Val}
  \quad\text{(variable env.)} \\[2pt]
\A     & : & \textit{AgentId} \rightharpoonup \textit{AgentDef}
  \quad\text{(agent env.)} \\[2pt]
\Gamma & : & \textit{ConfigKey} \rightharpoonup \textit{Val}
  \quad\text{(config env.)}
\end{array}
\]
\[
\textit{AgentDef} \;::=\; \backend(b)
  \mid \fixed(k)
  \mid \custom(\textit{prompt},m)
\]

In the Haskell implementation all three environments are bundled in a single
record:

\begin{lstlisting}[style=hs]
data EvalState = EvalState
  { sVar    :: Map String Value    -- sigma
  , sAgent  :: Map String AgentDef -- A
  , sConfig :: Map String Value    -- Gamma
  }
\end{lstlisting}

% =============================================================================
\section{Big-Step Operational Semantics}
\label{sec:semantics}

We define two evaluation relations following the big-step style of Plotkin~\cite{plotkin}:

\begin{itemize}
  \item $\J\; e \Downarrow v$ ---
        expression $e$ evaluates to value $v$
  \item $\J\; s \Downarrow \sigma';\Gamma'$ ---
        statement $s$ transforms $(\sigma,\Gamma)$ to $(\sigma',\Gamma')$
\end{itemize}

The agent environment $\A$ is updated mutably by declaration statements; it
is read-only during expression evaluation.

\textbf{Error Propagation (E-Prop).}  For any composite form, if a
constituent sub-evaluation yields $\err{v}$, the enclosing construct
immediately evaluates to $\err{v}$ without executing further sub-terms.
This is implemented via the \texttt{bindE} combinator:

\begin{lstlisting}[style=hs]
-- E-Prop: short-circuit on EErr
bindE :: EvalIO a -> (a -> EvalIO b)
      -> EvalIO b
bindE m f = do
  r <- m
  case r of
    EOk a  -> f a
    EErr v -> return (EErr v)
\end{lstlisting}

\subsection{Expression Rules}

\begin{mathpar}
\inferrule*[right=\textsc{E-Const}]
  { }
  { \J c \Downarrow c }

\inferrule*[right=\textsc{E-Var}]
  { \sigma(x) = v }
  { \J x \Downarrow v }

\inferrule*[right=\textsc{E-VarErr}]
  { x \notin \Dom(\sigma) }
  { \J x \Downarrow \err{\texttt{"unbound variable"}}}
\end{mathpar}

\begin{mathpar}
\inferrule*[right=\textsc{E-Proj}]
  { \J e \Downarrow \{\ldots,\,f{=}v,\,\ldots\} }
  { \J e.f \Downarrow v }
\end{mathpar}

\begin{mathpar}
\inferrule*[right=\textsc{E-Op}]
  { \J e_1 \Downarrow v_1 \\
    \J e_2 \Downarrow v_2 \\
    v = \sem{\oplus}(v_1,v_2) }
  { \J e_1 \oplus e_2 \Downarrow v }
\end{mathpar}

\begin{mathpar}
\inferrule*[right=\textsc{E-Rec}]
  { \J e_i \Downarrow v_i \quad \forall\,i \in 1..n }
  { \J \{f_i{=}e_i\} \Downarrow \{f_i{=}v_i\} }
\end{mathpar}

The three agent-dispatch rules correspond to the three \texttt{AgentDef}
constructors in \texttt{Syntax.hs}:

\begin{mathpar}
\inferrule*[right=\textsc{E-Agent-Backend}]
  { \A(A) = \backend(b) \\
    \J e_i \Downarrow v_i \quad \forall\,i \\
    b(\Gamma,v_1,\ldots,v_n) \Rightarrow v }
  { \J A(e_1,\ldots,e_n) \Downarrow v }
\end{mathpar}

\begin{mathpar}
\inferrule*[right=\textsc{E-Agent-Fixed}]
  { \A(A) = \fixed(k) \\
    \J e_i \Downarrow v_i \quad \forall\,i \\
    \runfixed(k,v_1,\ldots,v_n) \Rightarrow v }
  { \J A(e_1,\ldots,e_n) \Downarrow v }
\end{mathpar}

\begin{mathpar}
\inferrule*[right=\textsc{E-Agent-Custom}]
  { \A(A) = \custom(\textit{prompt},m) \\
    \J e \Downarrow v_{\textit{in}} \\
    \llmcall(\Gamma,m,\textit{prompt},v_{\textit{in}}) \Rightarrow v }
  { \J A(e) \Downarrow v }
\end{mathpar}

\begin{mathpar}
\inferrule*[right=\textsc{E-Pipe}]
  { \A(A) = d \\
    \J e \Downarrow v \\
    \apply(d,\Gamma,v) \Rightarrow v' }
  { \J e \rhd A \Downarrow v' }
\end{mathpar}

\subsection{Statement Rules}

\begin{mathpar}
\inferrule*[right=\textsc{S-Config}]
  { \J e_i \Downarrow v_i \quad \forall\,i }
  { \J \texttt{config}\;\{c_i{=}e_i\}
      \Downarrow \sigma;\,\Gamma[c_i{\mapsto}v_i] }
\end{mathpar}

The three declaration rules install the agent into $\A$; the variable and
config environments are passed through unchanged:

\begin{mathpar}
\inferrule*[right=\textsc{S-Decl-Backend}]
  { }
  { \J \texttt{agent}\;A\;\texttt{from}\;b \;\Downarrow\; \sigma;\Gamma \\\\
    (\A := \A[A \mapsto \backend(b)]) }
\end{mathpar}

\begin{mathpar}
\inferrule*[right=\textsc{S-Decl-Fixed}]
  { }
  { \J \texttt{agent}\;A{=}\texttt{FixedAgent}(k) \;\Downarrow\; \sigma;\Gamma \\\\
    (\A := \A[A \mapsto \fixed(k)]) }
\end{mathpar}

\begin{mathpar}
\inferrule*[right=\textsc{S-Decl-Custom}]
  { \J e \Downarrow v_p }
  { \J \texttt{agent}\;A{=}\texttt{CustomAI}(\ldots)
        \;\Downarrow\; \sigma;\Gamma \\\\
    (\A := \A[A \mapsto \custom(v_p,m)]) }
\end{mathpar}

\begin{mathpar}
\inferrule*[right=\textsc{S-Let}]
  { \J e \Downarrow v }
  { \J \texttt{let}\;x{=}e \;\Downarrow\; \sigma[x{\mapsto}v];\,\Gamma }

\inferrule*[right=\textsc{S-Seq}]
  { \J s_1 \Downarrow \sigma';\Gamma' \\
    \A;\sigma';\Gamma' \vdash s_2 \Downarrow \sigma'';\Gamma'' }
  { \J s_1\,;\,s_2 \;\Downarrow\; \sigma'';\Gamma'' }
\end{mathpar}

\begin{mathpar}
\inferrule*[right=\textsc{S-IfTrue}]
  { \J e \Downarrow \texttt{true} \\
    \J s_1 \Downarrow \sigma';\Gamma' }
  { \J \texttt{if}\;e\;\texttt{then}\;s_1\;\texttt{else}\;s_2
      \Downarrow \sigma';\Gamma' }
\end{mathpar}

\begin{mathpar}
\inferrule*[right=\textsc{S-IfFalse}]
  { \J e \Downarrow \texttt{false} \\
    \J s_2 \Downarrow \sigma';\Gamma' }
  { \J \texttt{if}\;e\;\texttt{then}\;s_1\;\texttt{else}\;s_2
      \Downarrow \sigma';\Gamma' }
\end{mathpar}

\begin{mathpar}
\inferrule*[right=\textsc{S-Fail}]
  { \J e \Downarrow v }
  { \J \texttt{fail}\;e \;\Downarrow\; \err{v} }
\end{mathpar}

\subsection{Extension Rules: Retry and Try/Catch}

The \texttt{retry} construct re-executes a body statement up to $n$
additional times on failure.  The three rules cover the three cases:

\begin{mathpar}
\inferrule*[right=\textsc{S-Retry-OK}]
  { \J s \Downarrow \sigma';\Gamma' }
  { \J \texttt{retry}\;n\;s \Downarrow \sigma';\Gamma' }
\end{mathpar}

\begin{mathpar}
\inferrule*[right=\textsc{S-Retry-Step}]
  { \J s \Downarrow \err{v} \quad n > 0 \\
    \J \texttt{retry}\;(n{-}1)\;s \Downarrow r }
  { \J \texttt{retry}\;n\;s \Downarrow r }
\end{mathpar}

\begin{mathpar}
\inferrule*[right=\textsc{S-Retry-Done}]
  { \J s \Downarrow \err{v} }
  { \J \texttt{retry}\;0\;s \Downarrow \err{v} }
\end{mathpar}

The \texttt{try/catch} construct intercepts an error value and makes it
available to the handler under a user-chosen variable name:

\begin{mathpar}
\inferrule*[right=\textsc{S-Try-OK}]
  { \J s_1 \Downarrow \sigma';\Gamma' }
  { \J \texttt{try}\;s_1\;\texttt{catch}\;x{\Rightarrow}s_2
      \Downarrow \sigma';\Gamma' }
\end{mathpar}

\begin{mathpar}
\inferrule*[right=\textsc{S-Try-Catch}]
  { \J s_1 \Downarrow \err{v} \\
    \A;\sigma[x{\mapsto}v];\Gamma \vdash s_2 \Downarrow \sigma';\Gamma' }
  { \J \texttt{try}\;s_1\;\texttt{catch}\;x{\Rightarrow}s_2
      \Downarrow \sigma';\Gamma' }
\end{mathpar}

% =============================================================================
\section{Implementation}
\label{sec:implementation}

The AWL toolchain is implemented in approximately 1{,}000 lines of Haskell
across four modules, built with Cabal.  Runtime dependencies are
\texttt{parsec}, \texttt{containers}, \texttt{aeson},
\texttt{http-conduit}, \texttt{bytestring}, \texttt{text}, and
\texttt{vector}.

\begin{center}
\small
\begin{tabular}{ll}
\toprule
Module & Role \\
\midrule
\texttt{Syntax.hs}       & Pure AST types \\
\texttt{Parser.hs}       & Lexer + Parsec parser \\
\texttt{Interpreter.hs}  & Big-step evaluator \\
\texttt{Main.hs}         & CLI driver \\
\bottomrule
\end{tabular}
\end{center}

\texttt{Syntax.hs} is intentionally free of I/O or Parsec imports so that
both the parser and interpreter can depend on it without circular
dependencies.

\subsection{Lexer and Tokenizer}

The lexer is built with Parsec's \texttt{Text.Parsec.Token} module.  A
\texttt{LanguageDef} record specifies the token structure of AWL:

\begin{lstlisting}[style=hs, caption={Lexer definition -- Parser.hs}]
langDef = emptyDef
  { PT.commentLine     = "//"
  , PT.commentStart    = "/*"
  , PT.commentEnd      = "*/"
  , PT.identStart      = letter <|> char '_'
  , PT.identLetter     = alphaNum <|> char '_'
  , PT.reservedOpNames =
      ["=","==","!=",">","<",">=","<="
      ,"+","-","*","/","&&","||",".","=>",":"]
  , PT.reservedNames   = reservedNamesList
  , PT.caseSensitive   = True }
\end{lstlisting}

Reserved words include \texttt{config}, \texttt{agent}, \texttt{from},
\texttt{let}, \texttt{if}, \texttt{then}, \texttt{else}, \texttt{fail},
\texttt{retry}, \texttt{try}, \texttt{catch}, \texttt{print},
\texttt{FixedAgent}, \texttt{CustomAI}, \texttt{true}, \texttt{false},
and \texttt{null}.  The call to \texttt{PT.makeTokenParser} derives
whitespace-skipping parsers for identifiers, string and numeric literals,
operators, and delimiters automatically.

\subsection{Parser}

The parser transforms a token stream into a single right-nested \texttt{SSeq}
tree.  Its public entry point is:

\begin{lstlisting}[style=hs]
parseProgram :: FilePath -> String
             -> Either ParseError Stmt
parseProgram = P.parse program
\end{lstlisting}

\noindent\textbf{Expression parsing} uses \texttt{buildExpressionParser} over
a five-level precedence table:

\begin{center}
\small
\begin{tabular}{lll}
\toprule
Level & Operators & Assoc. \\
\midrule
1 (high) & \texttt{* /}           & Left \\
2        & \texttt{+ -}           & Left \\
3        & \texttt{== != > < >= <=} & None \\
4        & \texttt{\&\&}          & Left \\
5 (low)  & \texttt{||}            & Left \\
\bottomrule
\end{tabular}
\end{center}

Postfix field projections (\texttt{e.f}) chain as a left-recursive postfix
rule applied after any atom:

\begin{lstlisting}[style=hs]
postfix = atom >>= go
  where
    go e = (do reservedOp "."
               f <- fieldName
               go (EProj e f))
           <|> return e
\end{lstlisting}

\noindent\textbf{Statement parsing} uses \texttt{P.choice} over eight
alternatives.  Braced blocks \texttt{\{ s$_1$; s$_2$; \ldots \}} desugar to
right-nested \texttt{SSeq} nodes, directly matching the S-Seq rule:

\begin{lstlisting}[style=hs]
stmtBlock = do
  void (symbol "{")
  ss <- sepEndBy1 stmt semi
  void (symbol "}")
  return (foldr1 SSeq ss)
\end{lstlisting}

Agent declarations branch on the token following the agent name:

\begin{lstlisting}[style=hs]
stmtAgent = do
  reserved "agent"; name <- identifier
  choice
    [ do reserved "from"; b <- backendP
         return (SAgentBackend name b)
    , do reservedOp "="; agentRhs name ]
\end{lstlisting}

\subsection{Abstract Syntax Tree}

The \texttt{Syntax} module defines all AST types.  The key types are:

\begin{lstlisting}[style=hs, caption={Value -- Syntax.hs}]
data Value
  = VString String | VNumber Double
  | VBool   Bool   | VList   [Value]
  | VRecord [(String, Value)]
  | VNull
  deriving (Eq, Show)
\end{lstlisting}

\begin{lstlisting}[style=hs, caption={Expr -- Syntax.hs}]
data Expr
  = EConst  Value             -- c
  | EVar    String            -- x
  | EProj   Expr String       -- e.f
  | ECall   String [Expr]     -- A(e1,...,en)
  | EBin    BinOp Expr Expr   -- e1 op e2
  | EList   [Expr]
  | ERecord [(String, Expr)]  -- {fi=ei}
  deriving (Eq, Show)
\end{lstlisting}

\begin{lstlisting}[style=hs, caption={Stmt -- Syntax.hs}]
data Stmt
  = SConfig      [(String, Expr)]
  | SAgentBackend String Backend
  | SAgentFixed   String Kind
  | SAgentCustom  String Expr (Maybe String)
  | SLet         String Expr
  | SIf          Expr Stmt Stmt
  | SSeq         Stmt Stmt
  | SFail        Expr
  | SRetry       Int  Stmt
  | STryCatch    Stmt String Stmt
  | SPrint       Expr
  deriving (Eq, Show)
\end{lstlisting}

\subsection{Interpreter / Evaluator}

The interpreter implements the big-step rules as a monadic evaluator in
\texttt{IO}.  The result type \texttt{EResult~a} directly mirrors the
semantic outcome distinction:

\begin{lstlisting}[style=hs]
data EResult a = EOk a | EErr Value
  deriving (Eq, Show)
type EvalIO a = IO (EResult a)
\end{lstlisting}

Each \texttt{evalExpr} case is labelled with its corresponding rule
name.  Variable lookup (E-Var / E-VarErr):

\begin{lstlisting}[style=hs]
-- (E-Var / E-VarErr)
evalExpr st (EVar x) =
  case Map.lookup x (sVar st) of
    Just v  -> okE v
    Nothing -> errE (VString
      ("unbound variable: " ++ x))
\end{lstlisting}

Field projection (E-Proj):

\begin{lstlisting}[style=hs]
-- (E-Proj)
evalExpr st (EProj e f) =
  evalExpr st e `bindE` \v ->
    case v of
      VRecord fs -> case lookup f fs of
        Just v' -> okE v'
        Nothing -> errE (VString
          ("no such field: " ++ f))
      _ -> errE (VString
             ("field access on non-record"))
\end{lstlisting}

Binary operator evaluation (E-Op):

\begin{lstlisting}[style=hs]
-- (E-Op)
evalExpr st (EBin op e1 e2) =
  evalExpr st e1 `bindE` \v1 ->
  evalExpr st e2 `bindE` \v2 ->
    case applyBinOp op v1 v2 of
      Just v  -> okE v
      Nothing -> errE (VString
        ("type error in: " ++ show op))
\end{lstlisting}

The \texttt{applyBinOp} helper implements the denotation $\sem{\oplus}$:

\begin{lstlisting}[style=hs]
applyBinOp OpAdd (VNumber a)(VNumber b)
  = Just (VNumber (a + b))
applyBinOp OpAdd (VString a)(VString b)
  = Just (VString (a ++ b))   -- concat
applyBinOp OpGt  (VNumber a)(VNumber b)
  = Just (VBool (a > b))
applyBinOp OpAnd (VBool a)  (VBool b)
  = Just (VBool (a && b))
applyBinOp _ _ _ = Nothing   -- type error
\end{lstlisting}

Each \texttt{evalStmt} case is also labelled.  Sequencing (S-Seq):

\begin{lstlisting}[style=hs]
-- (S-Seq)
evalStmt st (SSeq s1 s2) =
  evalStmt st s1 `bindE` \st' ->
    evalStmt st' s2
\end{lstlisting}

The \texttt{if} statement checks for a boolean condition (S-IfTrue /
S-IfFalse):

\begin{lstlisting}[style=hs]
-- (S-IfTrue / S-IfFalse)
evalStmt st (SIf c s1 s2) =
  evalExpr st c `bindE` \v ->
    case v of
      VBool True  -> evalStmt st s1
      VBool False -> evalStmt st s2
      _ -> errE (VString
             "if-condition is not a boolean")
\end{lstlisting}

Retry (S-Retry-OK / S-Retry-Step / S-Retry-Done):

\begin{lstlisting}[style=hs]
-- (S-Retry-OK / Step / Done)
evalStmt st (SRetry n s) = loop n
  where
    loop 0 = evalStmt st s  -- final attempt
    loop k = do
      r <- evalStmt st s
      case r of
        EOk st' -> okE st'
        EErr _  -> loop (k - 1)
\end{lstlisting}

Try/catch (S-Try-OK / S-Try-Catch):

\begin{lstlisting}[style=hs]
-- (S-Try-OK / S-Try-Catch)
evalStmt st (STryCatch s1 x s2) = do
  r <- evalStmt st s1
  case r of
    EOk st' -> okE st'
    EErr v  ->
      evalStmt st { sVar =
        Map.insert x v (sVar st) } s2
\end{lstlisting}

\subsection{Agent Backends}

Agent execution is dispatched by \texttt{applyAgent} based on the
\texttt{AgentDef} stored in $\A$:

\begin{lstlisting}[style=hs]
applyAgent st (ADBackend b)    vs =
  runBackend (sConfig st) b vs
applyAgent st (ADFixed k)      vs =
  runFixedWorkflow (sConfig st) k vs
applyAgent st (ADCustom p m)  [v] =
  runLLM (sConfig st) m p v
applyAgent _  (ADCustom _ _)   _  =
  errE (VString "CustomAI: wrong arity")
\end{lstlisting}

\noindent\textbf{Backend behaviour:}

\begin{itemize}
  \item \texttt{mock:"v"} --- returns the wired constant immediately.
  \item \texttt{python:"f"} / \texttt{http:"u"} --- logs the call
        and returns a synthetic record
        \texttt{\{backend, target, args, output\}}.
  \item \texttt{llm:"m"} --- routes to the Anthropic Claude Messages
        API when \texttt{real\_llm = true} and
        \texttt{ANTHROPIC\_API\_KEY} is set; otherwise returns a
        deterministic stub.
\end{itemize}

\noindent\textbf{Fixed agents.}  The eleven \texttt{Kind} values have
deterministic stub implementations in \texttt{runFixed}.  For example,
the \textit{Critic} agent scores input based on keyword presence:

\begin{lstlisting}[style=hs]
runFixedOne Critic v =
  let txt    = lowerStr (valueText v)
      hasNeg = any (`isInfixOf` txt)
                ["bad","wrong","fail","error"]
      score  = if hasNeg then 0.4 else 0.9
  in VRecord
       [ ("score",   VNumber score)
       , ("critique", VString
           (if hasNeg then "negative indicators"
                      else "no issues")) ]
\end{lstlisting}

\noindent\textbf{CustomAI and live LLM.}  When \texttt{real\_llm = true}
and the API key is available, \texttt{runLLM} constructs a raw HTTP
POST to the Anthropic Messages API endpoint (no Haskell SDK exists):

\begin{lstlisting}[style=hs]
callClaudeWithSystem apiKey model sys inp = do
  let body = A.object
        [ "model"      .= model
        , "max_tokens" .= (1024 :: Int)
        , "system"     .= sys
        , "messages"   .=
            [A.object ["role"    .= "user"
                      ,"content" .= inp]]
        ]
  initReq <- parseRequest
    "POST https://api.anthropic.com/v1/messages"
  let req = setRequestHeader "x-api-key"
              [BS.pack apiKey] $ ...
  result <- try (httpLBS req)
  ...
\end{lstlisting}

The response field \texttt{content[0].text} is extracted, then
\texttt{llmOutputValue} attempts to parse it as JSON; on failure it
wraps the raw text in a \texttt{VString}.

\subsection{CLI Driver}

\texttt{Main.hs} provides the \texttt{awl} executable.  It accepts an
optional file path (defaulting to stdin) and three output mode flags:

\begin{center}
\small
\begin{tabular}{ll}
\toprule
Flag & Behaviour \\
\midrule
(none)            & Print all config and variable bindings \\
\texttt{--quiet}        & Suppress output on success \\
\texttt{--result-only}  & Print \texttt{vars["result"]} as text \\
\texttt{--result-json}  & Print \texttt{vars["result"]} as JSON \\
\bottomrule
\end{tabular}
\end{center}

A non-zero exit code is returned on both parse errors and evaluation
failures, enabling shell-level error detection.

% =============================================================================
\section{Workflow Example}
\label{sec:workflow}

We trace execution of the customer support triage pipeline
(\texttt{customer\_support\_triage.awl}) to illustrate how the semantic
rules compose at runtime.

\begin{lstlisting}[style=awl,
    caption={Customer support triage (abridged)}]
config { real_llm = true,
         llm_model = "claude-opus-4-7" };

agent Guard    = FixedAgent(Guardrail);
agent Extract  = FixedAgent(Extractor);
agent Route    = FixedAgent(Router);
agent Draft    = FixedAgent(Writer);
agent Validate = FixedAgent(Validator);
agent Critique = FixedAgent(Critic);
agent Summarize = FixedAgent(Summarizer);

let ticket = "Customer charged twice...";
let safety = Guard(ticket);
let facts  = Extract(safety.output);
let route  = Route(facts.output);
let draft  = Draft({ route = route.output,
                     facts = facts.output });
let validation = Validate(draft.output);
let critique   = Critique(validation.output);
let summary    = Summarize({
  route = route.output,
  draft = draft.output });

let result = { route = route.output,
               drafted_response = draft.output,
               summary = summary.output }
\end{lstlisting}

\subsection{Execution Trace}

Each \texttt{let} statement fires a sequence of semantic rules:
(1)~\textsc{E-Agent-Fixed} dispatches to \texttt{runFixedWorkflow};
(2)~\textsc{E-Proj} accesses the \texttt{.output} field of the returned
record; (3)~\textsc{S-Let} stores the result into $\sigma$.
The program executes as a chain of \textsc{S-Seq} reductions, threading the
updated state through each step:

\begin{center}
\small
\begin{tabular}{ll}
\toprule
Variable & Value (summarized) \\
\midrule
\texttt{ticket}     & \texttt{"Customer charged twice..."} \\
\texttt{safety}     & \texttt{\{safe=true, payload=...\}} \\
\texttt{facts}      & \texttt{\{extracted="Customer",...\}} \\
\texttt{route}      & \texttt{\{route="commerce",...\}} \\
\texttt{draft}      & \texttt{\{draft="draft about:..."\}} \\
\texttt{validation} & \texttt{\{valid=true, payload=...\}} \\
\texttt{critique}   & \texttt{\{score=0.9,...\}} \\
\texttt{summary}    & \texttt{\{summary="...",key\_points=[...]\}}\\
\texttt{result}     & \texttt{\{route=..., drafted\_response=...\}}\\
\bottomrule
\end{tabular}
\end{center}

\subsection{Pipeline Data Flow}

The linear agent chain is formed entirely by \textsc{E-Proj} feeding
the \texttt{.output} field of one agent's record into the argument of
the next:

\[
\textit{ticket}
\to \text{Guard}
\to \text{Extract}
\to \text{Route}
\to \text{Draft}
\to \text{Validate}
\to \text{Critique}
\to \text{Summarize}
\]

\subsection{Error Recovery Constructs}

The \texttt{retry} and \texttt{try/catch} constructs can be used to
make any step in the pipeline fault-tolerant without restructuring
surrounding code.  A demonstrative excerpt from \texttt{example.awl}:

\begin{lstlisting}[style=awl]
// Retry Summarize up to 2 extra times on fail
if critique.score > 0.8 then
  let result = summary
else
  retry 2 { let result = Summarize(raw) };

// Structured error recovery
try {
  fail "demonstration of failure handling"
} catch err => {
  let recovered = err
}
\end{lstlisting}

Under \textsc{S-Retry-Step}, each retry re-enters the body with the
\emph{same} initial $\sigma$---retries are stateless with respect to
the outer environment.  Under \textsc{S-Try-Catch}, the error payload
\texttt{err(v)} is bound to variable \texttt{err} in the handler,
making the message directly accessible as a value.

\subsection{Cloud Architecture Consensus Example}

The most complex example,
\texttt{cloud\_architecture\_consensus.awl}, demonstrates parallel
CustomAI agent usage and multi-input merging via the \textit{Merger}
fixed agent:

\begin{lstlisting}[style=awl]
agent AWSArchitect  = CustomAI(prompt = "...");
agent AzureArchitect = CustomAI(prompt = "...");
agent CostReviewer  = CustomAI(prompt = "...");
agent Merge = FixedAgent(Merger);
agent Judge = CustomAI(prompt = "...");

let aws_p   = AWSArchitect(use_case);
let azure_p = AzureArchitect(use_case);
let cost    = CostReviewer({
  aws = aws_p.output, azure = azure_p.output });

let packet  = Merge(aws_p.output,
                    azure_p.output,
                    cost.output);
let decision = Judge(packet.output);
\end{lstlisting}

The \textit{Merger} agent accepts a variable number of arguments
(matched by the multi-arity \textsc{E-Agent-Fixed} rule) and returns a
single consolidated record, which the downstream \textit{Judge} agent
then receives as its sole input.

% =============================================================================
\section{Testing}
\label{sec:testing}

The project includes 49 integration tests in \texttt{tests/*.awl}, executed
by \texttt{tests/run\_tests.sh}.  The harness builds the \texttt{awl}
binary via Cabal, runs each test with \texttt{ANTHROPIC\_API\_KEY} unset
(forcing deterministic stubs), and checks output against inline
directives:

\begin{center}
\small
\begin{tabular}{ll}
\toprule
Directive & Assertion \\
\midrule
\texttt{// EXPECT\_OK}     & Exit 0, stdout has \texttt{[ok]} \\
\texttt{// EXPECT\_ERR: s} & Non-zero exit, stderr contains $s$ \\
\texttt{// CHECK: s}       & Stdout contains substring $s$ \\
\texttt{// NOT: s}         & Stdout does not contain $s$ \\
\bottomrule
\end{tabular}
\end{center}

\noindent Coverage spans:

\begin{itemize}
  \item Arithmetic, string concatenation, and comparison operators.
  \item Boolean logic and short-circuit error propagation.
  \item Record construction, projection, and structural equality.
  \item All error-propagation paths (RHS of \texttt{let}, operands
        of binary ops, record fields, agent arguments).
  \item \texttt{if}-true, \texttt{if}-false, and non-boolean
        condition errors.
  \item \texttt{fail}, \texttt{try/catch}, and nested error binding.
  \item \texttt{retry} under both success and repeated failure.
  \item All four backend types: \texttt{mock}, \texttt{python},
        \texttt{http}, and \texttt{llm} stub.
  \item All eleven fixed-agent kinds.
  \item \texttt{CustomAI} correct and incorrect arity.
  \item \texttt{config} evaluation and LLM model selection.
  \item \texttt{parse\_json} and \texttt{stringify} builtins.
  \item Sequential variable shadowing via repeated \texttt{let}.
  \item Reserved words used as record field names.
\end{itemize}

% =============================================================================
\section{Conclusion and Future Work}
\label{sec:conclusion}

We have presented AWL, a declarative domain-specific language for AI agent
workflow composition.  The language provides a clean surface syntax for
declaring agents with heterogeneous backends, composing them with
sequential bindings and control flow, and recovering from failures using
\texttt{retry} and \texttt{try/catch}.  We gave a complete big-step
operational semantics using three runtime environments, and demonstrated
that the Haskell implementation is in direct correspondence with the
semantic rules---every \texttt{evalExpr} and \texttt{evalStmt} case is
annotated with its rule name.  A 49-case integration test suite validates
this correspondence under deterministic stub execution.

\medskip
\noindent\textbf{Future work} includes:

\begin{enumerate}
  \item \textbf{Static type system.}  A polymorphic record type system
        would catch field-access errors and agent arity mismatches at
        compile time, enabling a formal type-safety proof.

  \item \textbf{Parallel execution.}  The current semantics is strictly
        sequential.  A \texttt{parallel \{ \ldots \}} construct with an
        appropriate join semantics would improve throughput for independent
        agent calls.

  \item \textbf{Distributed runtime.}  Agent backends are currently
        invoked in-process.  A networked runtime would allow agents to
        run on separate machines or in sandboxed containers.

  \item \textbf{Formal metatheory.}  A proof of progress and preservation
        (or a suitable dynamically-typed analogue) would formally guarantee
        that well-formed programs either terminate with a value or an error
        and never get stuck.

  \item \textbf{Streaming outputs.}  LLM backends currently return
        complete responses.  Integrating streaming APIs would allow
        downstream agents to begin processing before an upstream agent
        finishes.
\end{enumerate}

% =============================================================================
\begin{thebibliography}{9}

\bibitem{langchain}
H.~Chase et al.
\textit{LangChain: Building Applications with LLMs through Composability}.
GitHub, 2022.
\url{https://github.com/langchain-ai/langchain}

\bibitem{autogen}
Q.~Wu et al.
\textit{AutoGen: Enabling Next-Gen LLM Applications via Multi-Agent
Conversation}.
arXiv:2308.08155, 2023.

\bibitem{n8n}
J.~Oberhauser.
\textit{n8n: Fair-Code Workflow Automation}.
GitHub, 2019.
\url{https://github.com/n8n-io/n8n}

\bibitem{plotkin}
G.~Plotkin.
\textit{A Structural Approach to Operational Semantics}.
Technical Report DAIMI FN-19, Aarhus University, 1981.

\bibitem{parsec}
D.~Leijen and E.~Meijer.
\textit{Parsec: Direct Style Monadic Parser Combinators for the Real World}.
Technical Report UU-CS-2001-35, Utrecht University, 2001.

\bibitem{pierce}
B.~C.~Pierce.
\textit{Types and Programming Languages}.
MIT Press, 2002.

\end{thebibliography}

\end{multicols}
\end{document}
